\documentclass[11pt]{article}
\usepackage[a4paper,margin=26mm]{geometry}
\usepackage{amsmath,amssymb,amsthm,mathtools,hyperref}
\setlength{\emergencystretch}{2em}
\newtheorem{theorem}{Theorem}
\newtheorem{lemma}[theorem]{Lemma}
\newtheorem{corollary}[theorem]{Corollary}
\newcommand{\Q}{\mathbb Q}\newcommand{\Z}{\mathbb Z}
\newcommand{\A}{\mathcal A}
\newcommand{\Mat}{\operatorname{Mat}}
\newcommand{\Spec}{\operatorname{Spec}}
\title{One rational coordinate for a finite family of algebraic witnesses}
\author{\(\mathbb Z\mathbb Q\) Diophantine Project}
\date{23 November 2026}
\begin{document}\maketitle
\begin{abstract}
We give an explicit integer polynomial whose one extra rational coordinate
encodes the existence of a rational character of a finite commutative
algebra on which a prescribed guard is nonzero. The construction tolerates
arbitrary nilpotents and arbitrarily many zero terminal square roots.
It uses the characteristic polynomial of all differences of generator
values, coefficient strata, a separating element, and one determinant.
Thus independent finite-root encoders can be merged without identifying
any of their ordinary rational witnesses. The construction has a stated
finite degree bound and is also applicable after separately justified
sharing of ordinary witnesses.
\end{abstract}

\section{Statement}
Put $R=\Z[T_1,\ldots,T_b]$. Let $\A$ be a commutative unital $R$-algebra
free of rank $n\ge1$ as an $R$-module. Fix a basis, elements
$\theta_1,\ldots,\theta_m\in\A$ generating
$\A\otimes_{\mathbb Z}\mathbb Q$ as a $\mathbb Q[T_1,\ldots,T_b]$-algebra, and a guard
$g\in\A$. For a rational specialization $t\in\Q^b$, write
$\A_t=\A\otimes_{R,t}\Q$. A character means a unital algebra homomorphism.

\begin{theorem}\label{main}
An explicit polynomial $F\in R[U]$ can be constructed such that
\begin{equation}\label{equivalence}
 \exists u\in\Q\quad F(t,u)=0
 \quad\Longleftrightarrow\quad
 \exists\phi:\A_t\longrightarrow\Q\quad\phi(g)\ne0
\end{equation}
for every $t\in\Q^b$. Neither reducedness nor etaleness of $\A_t$ is
assumed. In particular, zero roots in any number of quadratic stages are
permitted. If $d$ bounds the total degree of the multiplication-matrix
entries of $g,g\theta_1,\ldots,g\theta_m$, set
\[
 \ell=m+1,\qquad J=\ell(n+1)^2.
\]
One may take
\begin{equation}\label{degree}
 \deg F\le
 2n(J+1)\bigl[1+d\{1+2J(1+\ell(J+1))\}\bigr].
\end{equation}
The bound measures degree in all base coordinates and $U$.
\end{theorem}

The guard may be a product. It then requires each of its factors to be
nonzero at the selected character, without imposing nonvanishing at
unselected characters. No variables beyond $U$ are used in this theorem.

\section{Spectra, including nonreduced algebras}

\begin{lemma}\label{spectrum}
Let $A$ be a finite-dimensional commutative $\Q$-algebra. The eigenvalues
of multiplication by $a\in A$, over $\overline{\Q}$, are the values
$\sigma(a)$ of its geometric characters, with positive multiplicities.
Consequently, for two copies of its multiplication matrix, the eigenvalues
of $M_a\otimes I-I\otimes M_a$ are all differences
$\sigma(a)-\tau(a)$, with multiplicities.
\end{lemma}
\begin{proof}
The commuting multiplication operators can be simultaneously triangularized
over an algebraically closed field. Each diagonal entry is additive,
multiplicative and sends the identity to one, and so is a character.
Alternatively, decompose $A\otimes\overline{\Q}$ as a product of local
Artinian algebras; on each factor multiplication by $a$ is its residue
value plus a nilpotent operator. The tensor difference of upper triangular
matrices has the stated diagonal entries. Neither argument requires that
the nilpotent operators vanish.
\end{proof}

\section{A completely specified polynomial}
Define
\[
 \gamma_0=g,\qquad \gamma_i=g\theta_i\quad(1\le i\le m),
 \qquad A_i=M_{\gamma_i}\in\Mat_n(R).
\]
Set $N=n+1$ and adjoin an extra zero eigenvalue:
\[
 \widetilde A_i=\begin{pmatrix}A_i&0\\0&0\end{pmatrix},\qquad
 D_i=\widetilde A_i\otimes I_N-I_N\otimes\widetilde A_i.
\]
The extra eigenvalue is used to separate every permitted character from
zero. Compute the single monic polynomial
\begin{equation}\label{differencepolynomial}
 \chi(Z)=\prod_{i=0}^{m}\det(ZI_{N^2}-D_i)
       =Z^J+a_1Z^{J-1}+\cdots+a_J,
 \qquad a_0=1.
\end{equation}
Its roots comprise every difference of two character values of a
$\gamma_i$, and every such value itself and its negative.

The following are polynomials in the base coordinates:
\begin{align}
 R_0&=2+N+\sum_{i=0}^{m}\sum_{v,w=1}^{n}(A_i)_{vw}^2
                         +\sum_{j=1}^{J}a_j^2,\label{Rbound}\\
 B_0&=NR_0,\qquad K=R_0(2B_0)^{J-1}.\label{Kbound}
\end{align}
For each $s=0,\ldots,J$, put
\begin{align}
 \Delta_s&=a_s^2,& C_s&=\Delta_s+(4\ell B_0+2)K,\label{Cs}\\
 M_s&=\Delta_s\sum_{i=0}^{m}
                   C_s^{i+1}\Delta_s^{m-i}A_i,&
 H_s(U)&=\det(I_n-UM_s).\label{Ms}
\end{align}
Finally define
\begin{equation}\label{Fformula}
 \boxed{\displaystyle
 F(T,U)=\prod_{s=0}^{J}
       \left(H_s(U)^2+\sum_{j=s+1}^{J}a_j^2\right).}
\end{equation}
All empty sums are zero. Every determinant, sum and product is finite,
and every entry belongs to $R$. This is an explicit coefficient
specification for an integer polynomial, even when expanding it would
be impractical.

\section{Why the element separates the permitted characters}
Fix a rational base specialization. Thus all matrix entries and coefficients
below are real rational numbers, although their eigenvalues may be complex.
Let $s$ be the largest index with $a_s\ne0$, taking $s=0$ if all
$a_1,\ldots,a_J$ vanish. Then
\[
 \chi(Z)=Z^{J-s}(Z^s+a_1Z^{s-1}+\cdots+a_s).
\]
Every eigenvalue of $\widetilde A_i$ has absolute value at most $B_0$:
each matrix entry has absolute value at most $R_0$, and the maximum row
sum bounds the spectral radius. Hence every root of $\chi$ has absolute
value at most $2B_0$.

If $\delta$ is a nonzero root of $\chi$, the product of its $s$ nonzero
roots, counted with multiplicity, has absolute value $|a_s|$. Therefore
\begin{equation}\label{lowerseparation}
 |\delta|\ge\frac{|a_s|}{(2B_0)^{s-1}}
       \ge\frac{a_s^2}{R_0(2B_0)^{J-1}}
       =\frac{\Delta_s}{K}=:\varepsilon.
\end{equation}
We used $|a_s|\le R_0$ and $2B_0>1$. When $s=0$ there are no nonzero
roots and the following separation statements are vacuous unless a
permitted character exists; such a character would itself give a
nonzero root through $\gamma_0$ and the extra zero eigenvalue. Indeed,
when $s=0$ all $\gamma_i$ have only zero character values; then
$H_0(U)=1$ and there is no permitted character.

Since $\Delta_s>0$, put
\[
             B=\frac{C_s}{\Delta_s}
               =1+\frac{4\ell B_0+2}{\varepsilon}>1.
\]
For any two different vectors in
\[
 \bigl\{(\sigma(\gamma_0),\ldots,\sigma(\gamma_m)):\sigma
                 \text{ a geometric character}\bigr\}\cup\{0\},
\]
let $h$ be their largest differing coordinate. Its difference has absolute
value at least $\varepsilon$, by \eqref{lowerseparation}; every coordinate
difference has absolute value at most $2B_0$. The leading term of their
weighted difference satisfies
\[
 \left|B^{h+1}\delta_h\right|\ge\varepsilon B^{h+1},
 \qquad
 \left|\sum_{i<h}B^{i+1}\delta_i\right|
        \le2\ell B_0B^h<\varepsilon B^{h+1}.
\]
Thus the weighted sums are distinct. The matrix in \eqref{Ms} is
multiplication by
\[
 \beta_s=\Delta_s^{m+2}\sum_{i=0}^{m}B^{i+1}\gamma_i.
\]
The scalar factor is nonzero and does not affect separation.

If $\sigma(g)=0$, every $\sigma(\gamma_i)$ is zero and
$\sigma(\beta_s)=0$. If $\sigma(g)\ne0$, the vector is nonzero and
$\sigma(\beta_s)\ne0$. Finally, two characters with nonzero guard have
different vectors unless they are the same character: recover their
generator values as $\sigma(\theta_i)=\sigma(\gamma_i)/\sigma(\gamma_0)$.
Consequently $\beta_s$ separates all permitted characters and sends
exactly the excluded characters to zero.

\section{Rationality and the coefficient strata}

\begin{proof}[Proof of Theorem~\ref{main}, equivalence]
For the actual last nonzero coefficient index $s$, suppose first that a
permitted character $\phi:\A_t\to\Q$ exists. Its value
$\phi(\beta_s)$ is a nonzero rational number. Choose
$u=1/\phi(\beta_s)$. Lemma~\ref{spectrum} gives $H_s(u)=0$;
all $a_j$ for $j>s$ vanish, so \eqref{Fformula} is zero.

Conversely, suppose $F(t,u)=0$ with $t,u$ rational. One factor of
\eqref{Fformula} is zero. It is a sum of rational squares, so its
$H_s(u)$ and all its tail coefficients vanish. If $a_s=0$, then
$\Delta_s=0$, $M_s=0$ and $H_s(u)=1$, a contradiction. Thus $s$ is
exactly the largest nonzero coefficient index. The determinant equality
also implies $u\ne0$ and supplies a geometric character $\sigma$ with
\[
                         \sigma(\beta_s)=1/u\in\Q^\times.
\]
This character has nonzero guard. Every Galois conjugate of $\sigma$
has the same $\beta_s$ value, because that value is rational. Separation
therefore makes every such conjugate the same character. In particular,
every generator value $\sigma(\theta_i)$ is fixed by the absolute
Galois group of $\Q$, hence belongs to $\Q$. The character takes values
in $\Q$ and has nonzero guard, as required.
\end{proof}

\begin{proof}[Proof of the degree bound]
The entries of $D_i$ have degree at most $d$, so
$\deg a_j\le jd\le Jd$. It follows that
\[
 \deg R_0\le2Jd,\quad
 \deg\Delta_s\le2Jd,\quad
 \deg K\le2J^2d,\quad
 \deg C_s\le2J(J+1)d.
\]
In the sum defining $M_s$, the largest bound comes from $i=m=\ell-1$.
Every entry of that matrix has degree at most
\[
 d_\beta=d+2Jd+2\ell J(J+1)d
          =d\{1+2J(1+\ell(J+1))\}.
\]
Hence $\deg H_s\le n(1+d_\beta)$. This bound also dominates the
degrees needed for the squared tail coefficients in \eqref{Fformula}.
There are $J+1$ factors, proving \eqref{degree}.
\end{proof}

\section{Clearing denominators without losing selected characters}
A useful input model is a triangular finite-root specification. At each
stage a root $\theta$ satisfies a monic polynomial with coefficients
rational in the ordinary base coordinates and preceding roots. Collect
its denominators in a common expression $D$, and include $D$ in the
nonzero guard at the selected character. If
\[
 \theta^e+\sum_{j=0}^{e-1}\frac{N_j}{D}\theta^j=0,
\]
replace the finite root by $\xi=D\theta$. The equation becomes the monic
relation
\[
                 \xi^e+\sum_{j=0}^{e-1}N_jD^{e-1-j}\xi^j=0.
\]
On the guarded locus this is reversible by $\theta=\xi/D$. Continue
in the triangular order, clearing the new denominators after each
substitution. A finite tower of monic polynomial extensions is free, with
basis the monomials whose exponent in each generator is below its stage
degree. The rank is the product of the stage degrees. All source guards
become a single polynomial guard after multiplying their numerators and
the recorded denominators. Thus Theorem~\ref{main} applies.

This procedure does not require a denominator to be nonzero at every
geometric character. The product guard selects the allowed characters.
Characters at which a cleared denominator vanishes remain in the finite
algebra and are rejected by the same separator. No extra ordinary
rational coordinate is introduced when a finite root is rescaled.

\section{Merging independent encoders}

\begin{corollary}\label{merge}
Suppose each of $k$ finite-root predicates has $b_i$ ordinary rational
base coordinates, a finite free algebra of rank $n_i$, and a polynomial
nonzero guard. Their conjunction is expressible with all ordinary base
coordinates and just one additional rational coordinate. Its finite
algebra can be taken to be the tensor product, of rank $\prod_i n_i$,
and its guard to be the product of the guards. No simultaneous
realization assertion beyond the original separate predicates is needed.
\end{corollary}
\begin{proof}
A character of the tensor product into $\Q$ is exactly a collection of
characters of its factors into $\Q$, agreeing on the rational base.
Its product guard is nonzero exactly when all original guards are nonzero.
Apply Theorem~\ref{main}. In particular, a tensor product such as
$\Q[\epsilon,\delta]/(\epsilon^2,\delta^2)$ presents no difficulty;
monogenicity of the full nonreduced algebra is never used.
\end{proof}

Applied to a source integer predicate with $m$ total rational coordinates,
comprising $m-1$ ordinary base coordinates and one final finite-root encoder, $k$ independent input
tests use $k(m-1)+1$ coordinates including the $k$ inputs. The mixed source
with eight integer unknowns and one nonnegative unknown then uses
\[
                    9(m-1)+1-1+3=9m-6.
\]
Here the nonnegative input is replaced by $s^2+t^2+u^2+u$, whose integral
values are exactly the nonnegative integers. For the ordinary-$abc$
interface first feed each finite-root source $X_i^2+1$. This expression is
positive and is integral exactly when $X_i$ is integral, so the external
zero-input factor is unnecessary. Thus the project's
25-, 27- and ordinary-abc 11-coordinate finite-root interfaces give
respective unrestricted-degree transfer counts
\[
                              219,\qquad237,\qquad93.
\]
These applications retain all the arithmetic premises of the relevant
integer predicates. They need a different determinant specification from
nine separate polynomials. The old numerical degree bound must not be
reused without checking \eqref{degree} for the new matrices.

Further identification of ordinary base coordinates requires a separate
simultaneous-completeness proof. Once such a proof is supplied, the same
encoder applies to the resulting joint triangular root algebra. For a
shared root one includes its monic equation once, so consistency of its
selected value across all blocks is imposed automatically.
\end{document}
